\subsection{Interpretation in terms of homogeneous variations of Hodge structure}
As we learned from one of the referees, the cominuscule representation $\jepPubBbb{E}$ of $G$ has been interpreted in the language of Hodge structures, first by Gross \cite{jep93Gro94} in the tube type case, then by Sheng and Zuo \cite{jep93SZ10} in general. Indeed, the representation $\jepPubBbb{E}$ defines a canonical homogeneous complex variation of Hodge structure ($\mathbb C$-VHS) $E_{CY}:=\jepPubOplus_{p+q=r_M}E^{p,q}$ on $M$ or on any quotient $\Lambda\backslash M$, for $\Lambda$ a discrete subgroup of $G_{\mathbb R}$. The Hodge bundle $E^{p,q}$ is just the bundle over $M=G_{\mathbb R}/K_{\mathbb R}$ associated with the $K_{\mathbb R}$-module $\jepPubBbb{E}_{r_M-p}$. Geometrically, the corresponding filtration is given by the successive osculating tangent spaces to the cone over the subvariety $M\subset\mathbb P\jepPubBbb{E}$. This $\mathbb C$-VHS is of particular interest because it is of Calabi-Yau type, meaning that $\dim E^{r_M,0}=1$ and $E^{r_M-1,1}\simeq E^{r_M,0}\otimes\Omega_M^1$. It is a real variation of Hodge structure ($\mathbb R$-VHS) of weight $r_M$ precisely when $M$ is of tube type. This is exactly the same thing as saying, as we did, that $\jepPubBbb{E}$ is self-dual in this case.

From this point of view, what we do here can be interpreted in the following way. Let us work on the projectivized tangent bundle $\mathbb PT_M$ of $M$. The notion of rank of an element of $\mathfrak m^-$ naturally extends to elements in the tangent bundle of $M$, and the bundle $\mathbb PT_M$ is the union of $r_M$ $G_{\mathbb R}$-orbits $\operatorname{Orb}_1,\ldots,\operatorname{Orb}_{r_M}$, where $\operatorname{Orb}_r:=\{[\eta]\in\mathbb PT_M\mid\operatorname{rk}(\eta)=r\}$. It is also the quotient $G_{\mathbb R}/(K_{\mathbb R}\cap Z_r)$, where $Z_r$ is the centralizer of a nilpotent element $y$ in $\mathfrak m^-$ of rank $r$. The sets $\cup_{s\leqslant r}\operatorname{Orb}_s$ are the characteristic varieties of Mok \cite{jep93Mok89}. Let us fix $1\leqslant r\leqslant r_M$ and call $\pi_M$ the projection $\mathbb PT_M\to M$. The definition of the weight filtration $\jepPubBbb{W}^r_\jepPubBullet$ of $\jepPubBbb{E}$ associated with a nilpotent element $y\in\mathfrak m^-$ or rank $r$ can also be extended to give an increasing holomorphic weight filtration on the restriction of the $\mathbb C$-VHS $\pi_M^\star E_{CY}$ to $\operatorname{Orb}_r$, compatible with the decreasing Hodge filtration $(\jepPubOplus_{p\geqslant k}\pi_M^\star E^{p,q})_k$, and this defines a homogeneous family of ``complex'' mixed Hodge structures on $\operatorname{Orb}_r$. When $M$ has tube type, $\pi_M^\star E_{CY}$ is an $\mathbb R$-VHS and we have a real mixed Hodge structure on each fiber of $\pi_M^\star E_{CY}$ above $\operatorname{Orb}_r$.
\begin{example}
We keep the notations of Example~2.42, assuming now that we are in the tube type case: the group $G_{\mathbb R}$ is $\operatorname{SU}(p,p)$ with maximal compact subgroup $K_{\mathbb R}=\operatorname S(\operatorname U(p)\times\operatorname U(p))$. Let $C$, resp. $D$, be the subspace of $\mathbb C^{2p}$ generated by $(e_1,\ldots,e_p)$, resp. $(e_{p+1},\ldots,e_{2p})$. In this case, the cominuscule representation is
\[
\jepPubBbb{E}=\wedge^p(C\oplus D)=\jepPubOplus_{m+n=p}\wedge^m C\otimes\wedge^n D.
\]
Whenever $V$ is a complex vector space, we denote by $\overline V$ the complex vector space $V$ with the $\mathbb C$-action defined by $\lambda\cdot v=\overline\lambda v$. Note that, with this convention, a real structure on $V$ is the same as an involutory $\mathbb C$-linear isomorphism $\overline V\to V$. As representations of $K_{\mathbb R}$, $\overline C\simeq C^\star$, $\overline D\simeq D^\star$, and $\wedge^{2p}(C\oplus D)$ is the trivial representation. We get a $K_{\mathbb R}$-equivariant isomorphism
\[
\overline{\wedge^m C\otimes\wedge^n D}\simeq\wedge^m\overline C\otimes\wedge^n\overline D\simeq\wedge^m C^\star\otimes\wedge^n D^\star\simeq(\wedge^m C\otimes\wedge^n D)^\star\simeq\wedge^{p-m}C\otimes\wedge^{p-n}D,
\]
where the last isomorphism is given by the perfect pairing
\[
(\wedge^m C\otimes\wedge^n D)\otimes(\wedge^{p-m}C\otimes\wedge^{p-n}D)\longrightarrow\wedge^p C\otimes\wedge^p D\simeq\wedge^{2p}(C\oplus D)\simeq\mathbb C.
\]
Therefore we get a real structure on $\jepPubBbb{E}$, invariant by $K_{\mathbb R}$, using the isomorphism
\[
\overline{\jepPubBbb{E}}=\jepPubOplus_{m+n=p}\overline{\wedge^m C\otimes\wedge^n D}\simeq\jepPubOplus_{m+n=p}\wedge^n C\otimes\wedge^m D=\jepPubBbb{E}.
\]
Since we have tautologically, for this real structure, $\overline{\wedge^m C\otimes\wedge^n D}=\wedge^n C\otimes\wedge^m D$, this defines an $\mathbb R$-VHS of weight $p=r_M$ on the symmetric space $M=G_{\mathbb R}/K_{\mathbb R}$. This is the homogeneous VHS defined in \cite{jep93Gro94}.

If now $y$ is an element in $\mathfrak m^-$ of rank $r$, we moreover have the decompositions $C=A\oplus N$ and $D=B\oplus I$, as in Example~2.42, where the restriction of the action of $y$ to $A$ gives an isomorphism $A\simeq I$. The ($p$-shifted) increasing weight filtration $\jepPubBbb{W}^r_\jepPubBullet$ of $\jepPubBbb{E}$ associated to the action of $y$ is given by
\[
\jepPubBbb{W}^r_s=\jepPubOplus_{\substack{i-\ell\leqslant s-p\\i+j+k+\ell=p}}\wedge^i A\otimes\wedge^j N\otimes\wedge^k B\otimes\wedge^\ell I,
\]
whereas the decreasing Hodge filtration $\jepPubBbb{F}^\jepPubBullet$ is given by
\[
\jepPubBbb{F}^t=\jepPubOplus_{\substack{i+j\geqslant t\\i+j+k+\ell=p}}\wedge^i A\otimes\wedge^j N\otimes\wedge^k B\otimes\wedge^\ell I.
\]
Therefore the graded parts are:
\[
\operatorname{Gr}^{\jepPubBbb{F}}_t\operatorname{Gr}^{\jepPubBbb{W}^r}_s\jepPubBbb{E}=\jepPubOplus_{\substack{i-\ell=s-p\\i+j=t\\i+j+k+\ell=p}}\wedge^i A\otimes\wedge^j N\otimes\wedge^k B\otimes\wedge^\ell I.
\]
The intersection $Z_r\cap K_{\mathbb R}$ of $K_{\mathbb R}$ with the centralizer $Z_r$ of $y$ is the group
\[
\{(a,b,c)\in\operatorname U(r)\times\operatorname U(p-r)\times\operatorname U(p-r)\mid\det(a^2bc)=1\}.
\]
As $(Z_r\cap K_{\mathbb R})$-modules, we have:
\begin{align*}
\overline{\wedge^i A\otimes\wedge^j N\otimes\wedge^k B\otimes\wedge^\ell I}&\simeq\wedge^i A^\star\otimes\wedge^j N^\star\otimes\wedge^k B^\star\otimes\wedge^\ell I^\star\\
&\simeq(\wedge^i A^\star\otimes\wedge^j N^\star\otimes\wedge^k B^\star\otimes\wedge^\ell I^\star)\otimes\wedge^{2p}(A\oplus N\oplus B\oplus I)\\
&\simeq\wedge^{r-i}A\otimes\wedge^{p-r-j}N\otimes\wedge^{p-r-k}B\otimes\wedge^{r-\ell}I\\
&\simeq\wedge^{r-\ell}A\otimes\wedge^{p-r-j}N\otimes\wedge^{p-r-k}B\otimes\wedge^{r-i}I\\
&\simeq\wedge^{i'}A\otimes\wedge^{j'}N\otimes\wedge^{k'}B\otimes\wedge^{\ell'}I
\end{align*}
with $i'=r-\ell$, $j'=p-r-j$, $k'=p-r-k$ and $\ell'=r-i$ (on the fourth line we used the isomorphism $I\simeq A$). Therefore, if $i-\ell=s-p$ and $i+j=t$, then $i'-\ell'=i-\ell=s-p$ and $i'+j'=p-j-\ell=s-(i+j)=s-t$. Hence the real structure on $\jepPubBbb{E}$ defined as above with these isomorphisms is such that
\[
\overline{\operatorname{Gr}^{\jepPubBbb{F}}_t\operatorname{Gr}^{\jepPubBbb{W}^r}_s\jepPubBbb{E}}=\operatorname{Gr}^{\jepPubBbb{F}}_{s-t}\operatorname{Gr}^{\jepPubBbb{W}^r}_s\jepPubBbb{E}
\]
and the Hodge filtration induces a real Hodge structure of weight $s$ on $\operatorname{Gr}^{\jepPubBbb{W}^r}_s\jepPubBbb{E}$, so that the two filtrations $\jepPubBbb{W}^r_\jepPubBullet$ and $\jepPubBbb{F}^\jepPubBullet$ indeed define a real mixed Hodge structure on $\jepPubBbb{E}$.
\end{example}
However, again from the Hodge structures point of view, our purpose here is a bit different. We would like to construct a canonical ``weak sub-$\mathbb R$-VHS'' of weight $r$ of the restriction of the $\mathbb C$-VHS $\pi_M^\star E_{CY}$ to $\operatorname{Orb}_r$. This is the role played by the submodule $\jepPubBbb{V}^r$. The meaning of ``weak $\mathbb R$-VHS'' is that if $V^r$ denotes the equivariant holomorphic bundle $\jepPubBbb{V}^r\otimes\jepPubScript{O}_{\operatorname{Orb}_r}$ on $\operatorname{Orb}_r$ with its Gauss-Manin connection $\nabla$, and if $F^\jepPubBullet$ denotes the Hodge filtration on $V^r$ defined by the filtration $\jepPubBbb{F}^\jepPubBullet$ of $\jepPubBbb{E}$, then $(V^r,F^\jepPubBullet)$ does have the symmetries of an $\mathbb R$-VHS by Proposition~2.40, but does not satisfy Griffiths transversality $\nabla(F^t)\subset F^{t-1}\otimes\Omega^1_{\operatorname{Orb}_r}$. Rather, $(V^r,F^\jepPubBullet)$ satisfies a weak transversality condition by Proposition~2.29. Indeed we have $\nabla(F^t\otimes\jepPubScript{O}_{\mathbb PT_M}(-1)\jepPubRestrict_{\operatorname{Orb}_r})\subset F^{t-1}$, where $\jepPubScript{O}_{\mathbb PT_M}(-1)$ is the tautological line bundle on $\mathbb PT_M$. This will allow us to construct \emph{leafwise} Higgs subsheaves (but not plain Higgs subsheaves) in Section~4.3.

